\scanpage{69}
\mspage{75}

Pour résumer

% slanted
\margin{i.e. Hot-asphériques}

\uline{Proposition 1} : $p_X$ et $q_X$ sont des $\deleted{\unsure{\mathbb D}}\mathrm{Hot}$-équivalences strictes, à fibres les $(X_{/x})^\circ$ et $ {}_{\deleted{\unsure{y}}x\backslash}X$ respectivement (donc à fibres contractiles)

\margin{(\unsure{donc aussi|Cela en sera} $\mathbb D$-équiv. strictes si $\unsure{\mathscr D|\mathscr P}(X)$ dans dom de $\mathbb D$)}

\deleted{\uline{Corollaire} On a une équivalence} \deleted{On voudrait}

\uline{Corollaire 1} Pour que $X$ soit $\mathbb D$-contractile, il f. et s. que $X^\circ$ le soit.

Plus généralement

\uline{Corollaire 2} Soit $f:X\longrightarrow Y$. Pour que $f$ soit une $\mathbb D$-équivalence, il f. et s. que $f^\circ$ le soit.

On utilise que le diagramme
\diagram{\begin{tikzcd}[ampersand replacement=\&]
\& \mathscr D(X) \arrow[dl,"p_X"'] \arrow[dr,"q_X"] \& \\
X \& \& X^\circ
\end{tikzcd}}
est fonctoriel en $X$. \unsure{On a donc|Ceci a donc}
\diagram{\begin{tikzcd}[ampersand replacement=\&,column sep=small]
\& \mathscr D(X) \arrow[rr,leftrightarrow,"{\mathscr D(f)}"] \arrow[dl,"p_X"'] \arrow[dd,"q_X"] \& \& \mathscr D(Y) \arrow[dl,"p_Y"] \arrow[dd,"q_Y"]\\
X \arrow[rr,"f",crossing over] \& \& Y \& \\
\& X^\circ \arrow[rr,"f^\circ"] \& \& Y^\circ
\end{tikzcd}}
les deux carrés commutatifs. Comme $p_X,p_Y$ sont des $\mathbb D$-équiv., $f$ est une $\mathbb D$-équiv. sss $\mathscr D(f)$ l'est, et de $=$ pour $f^\circ$.

Ces raisonnements supposent que si $X$ est dans le domaine de $\mathbb D$, $\mathscr D(X)$ aussi. C'est le cas si $\mathrm{Diag}$ est formé des cat. ordonnées finies.

\figure{\begin{tikzpicture}\draw (0,.35) -- (.17,.18) -- (.34,0);\draw[line width=.8pt] (.14,.23) -- (.21,.13);\end{tikzpicture}}
% archive sheet number: 68
